For $p=100$, Ridge achieved a minimum test-grid MSE of 4.5384; the 5-fold minimum-CV and one-SE choices gave test MSEs 4.5538 and 6.3267, while the corresponding 10-fold choices gave 4.5538 and 5.5740. Lasso achieved a minimum test-grid MSE of 4.4370; the 5-fold minimum-CV and one-SE choices gave test MSEs 4.6753 and 6.2272, while the corresponding 10-fold choices gave 4.6373 and 5.8349. For $p=200$, Ridge achieved a minimum test-grid MSE of 5.4811; the 5-fold minimum-CV and one-SE choices gave test MSEs 5.9133 and 5.7514, while the corresponding 10-fold choices gave 5.8159 and 5.7514. Lasso achieved a minimum test-grid MSE of 4.0099; the 5-fold minimum-CV and one-SE choices gave test MSEs 4.0238 and 5.2401, while the corresponding 10-fold choices gave 4.0493 and 4.9054. Increasing the dimension adds correlated predictors without changing the four variables in the generating mean. The comparisons reflect both increased estimation complexity and the realized samples. At $p=200$, the centered design has rank at most 199, making the unpenalized slope solution nonunique, while positive ridge penalties remain well defined. The Lasso fraction uses the minimum-norm least-squares reference. Changing the dimension also changes the random-number sequence, so these are not paired comparisons on identical observations.
